% GENERATED FILE. Edit canonical lessons, then run npm run generate:books.
% canonical-source-hash: fnv1a64:21b1841c96d49ade
% canonical-lessons: LA-C546-cavea, LA-C546-pulpitum, LA-C546-aulaeum, LA-C546-prologus, LA-C546-pantomimus

\chapter{At the theatre}
\label{ch:la-at-the-theatre}
\hlchaptermodality{546}

\begin{chapteropening}
\textbf{By the end of this chapter:} \emph{I can describe a Roman theatre.}

Complete the last lesson of chapter 546: I can describe a Roman theatre.
\end{chapteropening}

\section[cavea, caveae]{cavea, caveae — the seats of a theatre; a cage}

\label{lesson:LA-C546-cavea}



\begin{quote}
Before the new one: say the Latin for a band of musicians, then the Latin for harmony.
\end{quote}

\subsection*{You'll want to know: cavea, caveae}
\textbf{cavea, caveae} — \textquotedblleft{}the seats of a theatre; a cage\textquotedblright{}.

\begin{quote}
Tōta cavea plaudit.
\end{quote}

\begin{quote}
\emph{The whole audience applauds.}
\end{quote}

Feminine, first declension, from \emph{cavus} — hollow. The hollow bowl of seats, and so the audience; also a cage or den. English \emph{cage}.

\begin{cousinweb}
\textbf{cavea} is built on \emph{cavus} \textquotedblleft{}hollow\textquotedblright{} (Lewis \& Short). English \emph{cage} comes from it through Old French.
\end{cousinweb}

\begin{grammarlens}[title={What you've built}]
One more word for telling, explaining and describing at B1.
\end{grammarlens}

\subsection*{Your turn}
\begin{itemize}
\raggedright
  \item \emph{Say it:} \emph{cavea, caveae}
  \item \emph{Say it:} \emph{cavea, caveae}, once more
  \item \emph{Say it:} a sentence of your own with \emph{cavea, caveae}
  \item \emph{Recall:} say the Latin for a band of musicians, then the Latin for harmony, then say \emph{cavea, caveae} again
\end{itemize}

\begin{tcolorbox}[breakable,colback=teal!4,colframe=teal!35!black,title={Before you move on}]
What does cavea, caveae mean? (\textquotedblleft{}The seats of a theatre; a cage\textquotedblright{}.) Say it once more.
\end{tcolorbox}
\section[pulpitum, pulpitī]{pulpitum, pulpitī — a stage, platform}

\label{lesson:LA-C546-pulpitum}



\begin{quote}
Before the new one: say the Latin for harmony, then the Latin for the seats of a theatre.
\end{quote}

\subsection*{You'll want to know: pulpitum, pulpitī}
\textbf{pulpitum, pulpitī} — \textquotedblleft{}a stage, platform\textquotedblright{}.

\begin{quote}
Histriō in pulpitum ascendit.
\end{quote}

\begin{quote}
\emph{The actor climbs onto the stage.}
\end{quote}

Neuter, second declension. A wooden platform for actors or speakers. English \emph{pulpit}.

\begin{cousinweb}
\textbf{pulpitum} is a Latin word whose origin is not settled; Lewis \& Short gives none. English \emph{pulpit} comes from it.
\end{cousinweb}

\begin{grammarlens}[title={What you've built}]
One more word for telling, explaining and describing at B1.
\end{grammarlens}

\subsection*{Your turn}
\begin{itemize}
\raggedright
  \item \emph{Say it:} \emph{pulpitum, pulpitī}
  \item \emph{Say it:} \emph{pulpitum, pulpitī}, once more
  \item \emph{Say it:} a sentence of your own with \emph{pulpitum, pulpitī}
  \item \emph{Recall:} say the Latin for harmony, then the Latin for the seats of a theatre, then say \emph{pulpitum, pulpitī} again
\end{itemize}

\begin{tcolorbox}[breakable,colback=teal!4,colframe=teal!35!black,title={Before you move on}]
What does pulpitum, pulpitī mean? (\textquotedblleft{}A stage, platform\textquotedblright{}.) Say it once more.
\end{tcolorbox}
\section[aulaeum, aulaeī]{aulaeum, aulaeī — a theatre curtain}

\label{lesson:LA-C546-aulaeum}



\begin{quote}
Before the new one: say the Latin for the seats of a theatre, then the Latin for a stage.
\end{quote}

\subsection*{You'll want to know: aulaeum, aulaeī}
\textbf{aulaeum, aulaeī} — \textquotedblleft{}a theatre curtain\textquotedblright{}.

\begin{quote}
Aulaeum premitur, et fābula incipit.
\end{quote}

\begin{quote}
\emph{The curtain is lowered, and the play begins.}
\end{quote}

Neuter, second declension. Roman curtains worked upside down: they sank into a slot at the start of a play and rose at the end.

\begin{cousinweb}
\textbf{aulaeum} was borrowed from Greek \emph{aulaia} \textquotedblleft{}curtain, hanging\textquotedblright{} (Lewis \& Short; EtymDB).
\end{cousinweb}

\begin{grammarlens}[title={What you've built}]
One more word for telling, explaining and describing at B1.
\end{grammarlens}

\subsection*{Your turn}
\begin{itemize}
\raggedright
  \item \emph{Say it:} \emph{aulaeum, aulaeī}
  \item \emph{Say it:} \emph{aulaeum, aulaeī}, once more
  \item \emph{Say it:} a sentence of your own with \emph{aulaeum, aulaeī}
  \item \emph{Recall:} say the Latin for the seats of a theatre, then the Latin for a stage, then say \emph{aulaeum, aulaeī} again
\end{itemize}

\begin{tcolorbox}[breakable,colback=teal!4,colframe=teal!35!black,title={Before you move on}]
What does aulaeum, aulaeī mean? (\textquotedblleft{}A theatre curtain\textquotedblright{}.) Say it once more.
\end{tcolorbox}
\section[prologus, prologī]{prologus, prologī — a prologue; its speaker}

\label{lesson:LA-C546-prologus}



\begin{quote}
Before the new one: say the Latin for a stage, then the Latin for a theatre curtain.
\end{quote}

\subsection*{You'll want to know: prologus, prologī}
\textbf{prologus, prologī} — \textquotedblleft{}a prologue; its speaker\textquotedblright{}.

\begin{quote}
Prologus argūmentum fābulae nārrat.
\end{quote}

\begin{quote}
\emph{The prologue tells the plot of the play.}
\end{quote}

Masculine, second declension, from Greek. Both the speech before a comedy and the actor who gave it. English \emph{prologue}.

\begin{cousinweb}
\textbf{prologus} was borrowed from Greek \emph{prologos}, built on \emph{pro} \textquotedblleft{}before\textquotedblright{} and \emph{logos} \textquotedblleft{}speech\textquotedblright{} (EtymDB; Lewis \& Short). English \emph{prologue} comes from it.
\end{cousinweb}

\begin{grammarlens}[title={What you've built}]
One more word for telling, explaining and describing at B1.
\end{grammarlens}

\subsection*{Your turn}
\begin{itemize}
\raggedright
  \item \emph{Say it:} \emph{prologus, prologī}
  \item \emph{Say it:} \emph{prologus, prologī}, once more
  \item \emph{Say it:} a sentence of your own with \emph{prologus, prologī}
  \item \emph{Recall:} say the Latin for a stage, then the Latin for a theatre curtain, then say \emph{prologus, prologī} again
\end{itemize}

\begin{tcolorbox}[breakable,colback=teal!4,colframe=teal!35!black,title={Before you move on}]
What does prologus, prologī mean? (\textquotedblleft{}A prologue; its speaker\textquotedblright{}.) Say it once more.
\end{tcolorbox}
\section[pantomīmus, pantomīmī]{pantomīmus, pantomīmī — a mime dancer}

\label{lesson:LA-C546-pantomimus}



\begin{quote}
Before the new one: say the Latin for a theatre curtain, then the Latin for a prologue.
\end{quote}

\subsection*{You'll want to know: pantomīmus, pantomīmī}
\textbf{pantomīmus, pantomīmī} — \textquotedblleft{}a mime dancer\textquotedblright{}.

\begin{quote}
Pantomīmus sine verbīs fābulam saltat.
\end{quote}

\begin{quote}
\emph{The mime dancer dances a story without words.}
\end{quote}

Masculine, second declension, from Greek. A solo dancer who told myths with gestures; the best were stars.

\begin{cousinweb}
\textbf{pantomīmus} was borrowed from Greek \emph{pantomimos}, built on \emph{pan} \textquotedblleft{}all\textquotedblright{} and \emph{mimeisthai} \textquotedblleft{}to imitate\textquotedblright{} (EtymDB; Lewis \& Short). English \emph{pantomime} comes from it.
\end{cousinweb}

\begin{grammarlens}[title={What you've built}]
One more word for telling, explaining and describing at B1.
\end{grammarlens}

\subsection*{Your turn}
\begin{itemize}
\raggedright
  \item \emph{Say it:} \emph{pantomīmus, pantomīmī}
  \item \emph{Say it:} \emph{pantomīmus, pantomīmī}, once more
  \item \emph{Say it:} a sentence of your own with \emph{pantomīmus, pantomīmī}
  \item \emph{Recall:} say the Latin for a theatre curtain, then the Latin for a prologue, then say \emph{pantomīmus, pantomīmī} again
\end{itemize}

\begin{tcolorbox}[breakable,colback=teal!4,colframe=teal!35!black,title={Before you move on}]
What does pantomīmus, pantomīmī mean? (\textquotedblleft{}A mime dancer\textquotedblright{}.) Say it once more.
\end{tcolorbox}
